% =============================================================================
%  Notas de clase — Sesión 08: Práctica competitiva con bits
%  Programación Avanzada — Universidad Panamericana
%  Compilar: pdflatex sesion08_practica_competitiva.tex
% =============================================================================
\documentclass[12pt, oneside]{article}
\usepackage{progavanzada}
\usepackage{array}
\usepackage{multirow}

\begin{document}

\portada{Sesión 08}{Práctica competitiva con bits}{2026}

\tableofcontents
\newpage

% =============================================================================
\section{Instrucciones de la sesión}
% =============================================================================

Esta sesión es trabajo en \textbf{equipo de tres personas}.

La idea es la misma que en cualquier maratón de programación: leer un problema,
entenderlo, encontrar el patrón matemático y verificarlo con los ejemplos.
En la siguiente sesión lo implementan.  Hoy, todo se resuelve \textbf{a mano}.

\begin{advertencia}
  \textbf{Entregable:} un documento (PDF generado de fotos o escaneado) con:
  \begin{enumerate}
    \item Los casos de prueba resueltos paso a paso para cada problema.
    \item Un algoritmo claro (pseudocódigo, diagrama de flujo o fórmula
          matemática) que funcione para cualquier entrada, no solo los ejemplos.
  \end{enumerate}
  Un algoritmo que diga ``veo el patrón'' no cuenta: debe poder ejecutarse
  mecánicamente paso a paso.
\end{advertencia}

\begin{consejo}
  \textbf{Estrategia recomendada para cada problema:}
  \begin{enumerate}
    \item Convertir los números a binario y observar qué pasa bit a bit.
    \item Resolver los ejemplos del enunciado a mano para verificar que
          entiendes el problema.
    \item Resolver los casos de práctica adicionales.
    \item Generalizar: ¿qué fórmula o proceso aplica para \textit{cualquier}
          número?
  \end{enumerate}
\end{consejo}

\newpage

% =============================================================================
\section{Problema 1: Los bits de la clase de programación}
% =============================================================================

\subsection{Enunciado}

Uno de tus compañeros inventó un algoritmo que opera sobre una cadena de bits
$A$ y va construyendo otra cadena $S$ (inicialmente vacía).  El proceso se
repite hasta que $A$ queda vacía:

\begin{itemize}
  \item Si $A$ tiene \textbf{longitud impar}: el bit que está exactamente en el
        centro se retira de $A$ y se agrega al final de $S$.
  \item Si $A$ tiene \textbf{longitud par}: los dos bits del centro se comparan.
        El \textbf{mayor} (o cualquiera en caso de empate) se agrega \textbf{primero}
        al final de $S$, seguido del \textbf{menor}; ambos se retiran de $A$.
\end{itemize}

Al terminar, se imprime la representación decimal de $S$ como número binario.
Si el resultado es mayor que $10^9+7$, imprime el residuo módulo $10^9+7$.

\textbf{Entrada:} primera línea con $k$ (número de casos); luego $k$ cadenas
binarias.

\textbf{Salida:} para cada caso, \texttt{Caso \#x: y}.

\subsection{Ejemplo del enunciado}

\begin{center}
\begin{tabular}{ll}
  \toprule
  \textbf{Entrada} & \textbf{Salida} \\
  \midrule
  \texttt{00000} & \texttt{Caso \#1: 0} \\
  \texttt{0101101} & \texttt{Caso \#2: 106} \\
  \texttt{100} & \texttt{Caso \#3: 2} \\
  \bottomrule
\end{tabular}
\end{center}

\textbf{Traza de \texttt{0101101}:}

\begin{center}
\begin{tabular}{clp{4cm}l}
  \toprule
  \textbf{Paso} & \textbf{A} & \textbf{Acción} & \textbf{S acumulado} \\
  \midrule
  1 & \texttt{0101101} (7, impar)  & extrae centro [3]\,=\,\texttt{1}          & \texttt{1} \\
  2 & \texttt{010101}  (6, par)    & centro [2,3]\,=\,\texttt{0,1}; mayor\,=\,\texttt{1}, menor\,=\,\texttt{0} & \texttt{110} \\
  3 & \texttt{0101}    (4, par)    & centro [1,2]\,=\,\texttt{1,0}; mayor\,=\,\texttt{1}, menor\,=\,\texttt{0} & \texttt{11010} \\
  4 & \texttt{01}      (2, par)    & centro [0,1]\,=\,\texttt{0,1}; mayor\,=\,\texttt{1}, menor\,=\,\texttt{0} & \texttt{1101010} \\
  \bottomrule
\end{tabular}
\end{center}

$S = \texttt{1101010}_2 = 106_{10}$ \checkmark

\subsection{Casos de práctica}

Resuelve los siguientes casos a mano.  Llena la tabla de traza y escribe el
resultado final.

\begin{center}
\begin{tabular}{ccc}
  \toprule
  \textbf{A} & \textbf{S (binario)} & \textbf{Resultado decimal} \\
  \midrule
  \texttt{1}    & & \\[8pt]
  \texttt{110}  & & \\[8pt]
  \texttt{1010} & & \\[8pt]
  \texttt{1111} & & \\[8pt]
  \bottomrule
\end{tabular}
\end{center}

\vspace{1cm}
\textbf{Espacio para la traza:}

\vspace{4cm}

\begin{consejo}
  Convierte la cadena a posiciones numeradas de \textbf{izquierda a derecha}
  empezando en 0.  El índice del centro izquierdo en una cadena de longitud $n$
  par es $n/2 - 1$; el centro derecho es $n/2$.
\end{consejo}

\newpage

% =============================================================================
\section{Problema 2: Chasquidos Mágicos}
% =============================================================================

\subsection{Enunciado}

Tienes $n$ enchufes inteligentes conectados en serie, todos inicialmente
\textbf{apagados}.  Un enchufe cambia de estado al escuchar un chasquido,
pero \textbf{solo si recibe corriente}.  La corriente llega al primer enchufe
desde la pared; cada enchufe subsiguiente recibe corriente solo si el anterior
está \textbf{encendido}.  La lámpara está conectada al último enchufe ($n$-ésimo).

Dados $n$ (número de enchufes, $1 \le n \le 30$) y $c$ (número de chasquidos,
$0 \le c \le 10^8$), determina si la lámpara quedó \textbf{PRENDIDA} o
\textbf{APAGADA}.

\subsection{Ejemplo del enunciado}

\begin{center}
\begin{tabular}{ccp{6cm}}
  \toprule
  $n$ & $c$ & \textbf{Resultado} \\
  \midrule
  1 & 0  & \texttt{APAGADA} \\
  1 & 1  & \texttt{PRENDIDA} \\
  4 & 0  & \texttt{APAGADA} \\
  4 & 47 & \texttt{PRENDIDA} \\
  \bottomrule
\end{tabular}
\end{center}

\subsection{Casos de práctica}

Traza la cadena de enchufes para \textbf{n=2} y \textbf{n=3}.
Escribe el estado de cada enchufe (0=apagado, 1=encendido) después de cada
chasquido.  El enchufe $n$ determina si la lámpara está prendida.

\textbf{n = 2} (enchufe 2 = lámpara):

\begin{center}
\begin{tabular}{c|c|c|c}
  \toprule
  \textbf{Chasquidos} & \textbf{E1} & \textbf{E2} & \textbf{Lámpara} \\
  \midrule
  0 & 0 & 0 & APAGADA \\
  1 & & & \\[4pt]
  2 & & & \\[4pt]
  3 & & & \\[4pt]
  4 & & & \\[4pt]
  \bottomrule
\end{tabular}
\end{center}

\vspace{0.5cm}

\textbf{n = 3} (enchufe 3 = lámpara):

\begin{center}
\begin{tabular}{c|c|c|c|c}
  \toprule
  \textbf{Chasquidos} & \textbf{E1} & \textbf{E2} & \textbf{E3} & \textbf{Lámpara} \\
  \midrule
  0 & 0 & 0 & 0 & APAGADA \\
  1 & & & & \\[4pt]
  2 & & & & \\[4pt]
  3 & & & & \\[4pt]
  4 & & & & \\[4pt]
  5 & & & & \\[4pt]
  6 & & & & \\[4pt]
  7 & & & & \\[4pt]
  8 & & & & \\[4pt]
  \bottomrule
\end{tabular}
\end{center}

\vspace{0.5cm}
¿Qué patrón observas entre la columna ``Chasquidos'' escrita en binario y el
estado de cada enchufe?

\vspace{1.5cm}

\begin{consejo}
  Después de llenar las tablas, escribe $c$ en binario al lado de cada fila y
  compara columna por columna.  La respuesta se vuelve muy obvia.
\end{consejo}

\newpage

% =============================================================================
\section{Problema 3: Separación de bits — posiciones pares e impares}
% =============================================================================

\subsection{Enunciado}

Dado un número entero $n$ (menor a $2^{31}-1$) en base 10, separa sus bits en
dos números $a$ y $b$:

\begin{itemize}
  \item Los bits de $n$ que están en \textbf{posición impar} (bit 1, bit 3,
        bit 5, \ldots, contando desde 0 en el bit menos significativo) y que
        valen 1, forman el número $a$.  Todos los demás bits de $a$ son 0.
  \item Los bits de $n$ que están en \textbf{posición par} (bit 0, bit 2,
        bit 4, \ldots) y que valen 1, forman el número $b$.
\end{itemize}

Leer enteros hasta recibir un 0.  Por cada número imprimir ``$a$ $b$'' en una
línea.

\subsection{Ejemplo del enunciado}

\begin{center}
\begin{tabular}{ccccc}
  \toprule
  $n$ & \textbf{Binario} & \textbf{Bits impares} ($a$) & \textbf{Bits pares} ($b$) & \textbf{Salida} \\
  \midrule
  13 & \texttt{1101} & \texttt{1000}=8  & \texttt{0101}=5 & \texttt{8 5} \\
   6 & \texttt{0110} & \texttt{0010}=2  & \texttt{0100}=4 & \texttt{2 4} \\
   7 & \texttt{0111} & \texttt{0010}=2  & \texttt{0101}=5 & \texttt{2 5} \\
  \bottomrule
\end{tabular}
\end{center}

\begin{recuerda}
  \textbf{Posiciones de bits:} el bit menos significativo (el que vale $2^0$)
  está en la posición 0 (par).  El siguiente (valor $2^1$) está en posición 1
  (impar).  Y así sucesivamente.
\end{recuerda}

\subsection{Casos de práctica}

Completa la siguiente tabla. Pista: escribe el número en binario y subraya
los bits en posición impar para obtener $a$, y los bits en posición par para $b$.

\begin{center}
\begin{tabular}{cccccc}
  \toprule
  $n$ & \textbf{Binario} & \textbf{Pos. impares (a)} & $a$ & \textbf{Pos. pares (b)} & $b$ \\
  \midrule
   3  & & & & & \\[6pt]
   5  & & & & & \\[6pt]
  10  & & & & & \\[6pt]
  12  & & & & & \\[6pt]
  \bottomrule
\end{tabular}
\end{center}

\vspace{0.5cm}

¿Hay una máscara binaria que, aplicada con AND a $n$, extraiga exactamente los
bits en posición impar?  ¿Y para los bits en posición par?

\vspace{2cm}

\begin{consejo}
  Escribe en binario la constante \texttt{0xAAAAAAAA} y la constante
  \texttt{0x55555555}.  Fíjate qué posiciones de bits están encendidas en cada
  una.
\end{consejo}

\newpage

% =============================================================================
\section{Problema 4: Separación de bits — enumerar 1s}
% =============================================================================

\subsection{Enunciado}

Dado $n$, se enumeran los bits que valen 1 en $n$ de \textbf{derecha a izquierda}
empezando a contar desde 1.  Los 1s enumerados con un número \textbf{impar}
(1.°, 3.°, 5.°, \ldots) forman el número $a$ (conservando su posición de bit
dentro de $n$).  Los 1s enumerados con un número \textbf{par} (2.°, 4.°, 6.°,
\ldots) forman el número $b$.

Leer enteros hasta recibir un 0.  Por cada número imprimir ``$a$ $b$''.

\subsection{Ejemplo del enunciado}

\begin{center}
\begin{tabular}{p{1.2cm}p{2cm}p{4.8cm}cc}
  \toprule
  $n$ & \textbf{Binario} & \textbf{Enumeración de 1s (der.\ $\to$ izq.)} & $a$ & $b$ \\
  \midrule
  13 & \texttt{1101}
     & bit0=1 (1°, impar), bit2=1 (2°, par), bit3=1 (3°, impar)
     & 9 & 4 \\
   6 & \texttt{0110}
     & bit1=1 (1°, impar), bit2=1 (2°, par)
     & 2 & 4 \\
   7 & \texttt{0111}
     & bit0=1 (1°, impar), bit1=1 (2°, par), bit2=1 (3°, impar)
     & 5 & 2 \\
  \bottomrule
\end{tabular}
\end{center}

\textbf{Verificación para n=13 (\texttt{1101}):}
El 1.° y 3.° uno están en bits 0 y 3 → $a = 2^0 + 2^3 = 1+8 = 9$.
El 2.° uno está en bit 2 → $b = 2^2 = 4$.

\subsection{Casos de práctica}

\begin{center}
\begin{tabular}{p{0.8cm}p{2cm}p{4.5cm}cc}
  \toprule
  $n$ & \textbf{Binario} & \textbf{Enumeración de 1s} & $a$ & $b$ \\
  \midrule
   3  & & & & \\[14pt]
   5  & & & & \\[14pt]
  10  & & & & \\[14pt]
  12  & & & & \\[14pt]
  \bottomrule
\end{tabular}
\end{center}

\vspace{0.5cm}

¿Observas alguna relación entre este problema y el anterior (posiciones pares
e impares)?  ¿La separación que obtienes aquí siempre coincide con la de
posiciones pares/impares?  ¿Cuándo sí y cuándo no?

\vspace{2cm}

\begin{consejo}
  Empieza por $n$s pequeños con muy pocos 1s, como $n = 2^k$ (potencias de 2).
  Observa qué ocurre.  Luego prueba con $n=3$ (dos 1s) y $n=7$ (tres 1s) donde
  el patrón se vuelve más claro.
\end{consejo}

% =============================================================================
\section{Resumen y entregable}
% =============================================================================

\begin{center}
\begin{tabular}{lp{7cm}}
  \toprule
  \textbf{Problema} & \textbf{Clave a encontrar} \\
  \midrule
  Algoritmo de Phil    & Proceso de extracción de bits centrales; traza
                         paso a paso del algoritmo \\[4pt]
  Chasquidos Mágicos   & Relación entre el estado de los enchufes y la
                         representación binaria de $c$ \\[4pt]
  Separación pares/imp & Máscara AND que aísla los bits en posiciones pares
                         o impares \\[4pt]
  Separación enum.\ 1s & Enumeración de los bits encendidos y distribución
                         entre $a$ y $b$ \\
  \bottomrule
\end{tabular}
\end{center}

\vspace{0.5cm}

\begin{advertencia}
  \textbf{Entrega:} sube a Canvas un PDF con fotos de:
  \begin{enumerate}
    \item Los casos de prueba de cada problema resueltos a mano.
    \item El algoritmo (pseudocódigo, fórmula o diagrama) para cada problema.
  \end{enumerate}
  Incluye el \textbf{nombre completo de los tres integrantes} del equipo en
  la primera página.
\end{advertencia}

\end{document}
